\section{Conclusiones}

Las decisiones de este proyecto estuvieron condicionadas desde el principio por el emplazamiento. La combinación de suelo desértico salino, alta evaporación, radiación intensa y zona remota no es un contexto neutral: determina el tipo de cemento, la duración del curado, la logística del suministro de hormigón y hasta la conveniencia de prefabricar en vez de construir en terreno. Eso hizo que varias especificaciones que podrían parecer conservadoras tengan una justificación directa en las condiciones del sitio.

Para el clarificador secundario, la geometría circular y la exigencia de impermeabilidad llevaron a proponer encofrado deslizante con waterstop en la junta losa--muro. El HAC se identificó como la opción adecuada para el muro no por ser innovador sino porque la densidad de armadura hace que la compactación mecánica convencional sea poco confiable en esa sección. El sobrecosto del 11\% respecto al G35 convencional es marginal frente al riesgo de una falla hidráulica en una estructura de contención de aguas residuales. Los hormigones G35, G25 y G17 se especificaron con exposición S2 dado el suelo desértico salino del sector, con cemento resistente a sulfatos para las estructuras en contacto con el terreno.

Los tres problemas de calidad analizados (nidos de piedras en el muro, fisuración por retracción plástica en las losas y filtración por junta fría) tienen en común que su costo de reparación supera con creces el de prevenirlos. La fisuración plástica es el riesgo más probable en La Negra porque la evaporación superficial es alta prácticamente durante todo el año, y no solo en verano, condición que el diseño de las medidas preventivas tuvo que considerar de forma permanente.

El campamento se resolvió con módulos estructurales 3D prefabricados, que reducen la mano de obra en terreno entre un 70 y un 90\% y eliminan la dependencia de las condiciones climáticas del desierto durante la fabricación. Para una obra en una zona de difícil acceso, eso se traduce directamente en menores costos de logística y menor exposición del personal.
